% Figure 2.2 --- Vietoris--Rips filtration and persistence barcode on a toy point cloud
\begin{figure}[!htb]
\centering
\resizebox{\linewidth}{!}{%
\begin{tikzpicture}[font=\footnotesize]
\def\pts{(0,0),(1.2,0.1),(1.6,1.1),(0.8,1.8),(-0.2,1.1),(3.0,0.9)}
\foreach \k/\r/\lab in {0/0.08/{$\varepsilon=0$},1/0.62/{$\varepsilon_1$},2/0.8/{$\varepsilon_2$},3/1.05/{$\varepsilon_3$}}{
 \begin{scope}[xshift=\k*4.3cm]
  \foreach \p in \pts { \fill[axisSf, opacity=0.55] \p circle (\r); }
  \ifnum\k>0
   \draw[axisS, thick] (0,0)--(1.2,0.1)--(1.6,1.1)--(0.8,1.8)--(-0.2,1.1)--cycle;
  \fi
  \ifnum\k>2
   \fill[axisS, opacity=0.35] (0,0)--(1.2,0.1)--(1.6,1.1)--(0.8,1.8)--(-0.2,1.1)--cycle;
   \draw[axisS, thick] (1.6,1.1)--(3.0,0.9);
  \fi
  \foreach \p in \pts { \fill[black] \p circle (1.6pt); }
  \node at (1.4,-1.3) {\lab};
 \end{scope}}
% barcode
\begin{scope}[xshift=0cm, yshift=-4.6cm]
 \draw[-{Latex[length=2mm]}] (0,0) -- (16.5,0) node[right] {scale $\varepsilon$};
 \foreach \x/\l in {0/0, 4.3/{\varepsilon_1}, 8.6/{\varepsilon_2}, 12.9/{\varepsilon_3}} \draw (\x,0.08)--(\x,-0.08) node[below] {$\l$};
 \node[left] at (-0.2,1.55) {$H_0$};
 \foreach \y/\e in {1.1/4.3,1.3/4.3,1.5/4.3,1.7/4.3} \draw[axisR, line width=2pt] (0,\y) -- (\e,\y);
 \draw[axisR, line width=2pt] (0,1.9) -- (12.9,1.9);
 \draw[axisR, line width=2pt, -{Latex[length=2mm]}] (0,2.1) -- (16.3,2.1);
 \node[left] at (-0.2,0.55) {$H_1$};
 \draw[failred, line width=2.5pt] (4.3,0.55) -- (11.5,0.55);
 \node[right, font=\scriptsize, text=failred] at (11.6,0.55) {one persistent loop};
\end{scope}
\end{tikzpicture}}
\caption[Vietoris--Rips filtration and persistence barcode]{Vietoris--Rips filtration on a
toy point cloud of six points and the resulting barcode. As the scale $\varepsilon$ grows,
balls around the points merge: components ($H_0$ bars) die when they join, a loop ($H_1$
bar, red) is born when the five points on the left close a cycle and dies when the cycle is
filled. Long bars correspond to robust structure, short bars to noise. In this thesis the
points are the branch points of the vessel skeleton.}
\label{fig:filtration}
\end{figure}
